\subsection{Integral Domains}

\begin{definition}
    A commutative ring with 1 and no zero divisors is
    called an integral domain
\end{definition}

\begin{proposition}
    If $R$ is an integral domain, $a,b,c\in R$, $a\neq 0$, and $ab=ac$
    then $b=c$
\end{proposition}

\begin{proof}
    $ab=ac \implies ab-ac=0 \implies a(b-c)=0$ since there are no zero
    divisors and $a\neq 0$ 
    
    $\implies (b-c)=0 \implies b=c$
\end{proof}

\begin{theorem}
    \label[theorem]{thm:integralfield}
    Every finite integral domain is a field
\end{theorem}

\begin{proof}
    Suppose $R$ is a finite integral domain.

    Let $0\neq r\in R$

    Consider: $A=\{r^k \mid k \geq 0\}=\{1,r,r^2,\dots\}\subseteq R$

    Since $A$ is finite $\exists m > n$ s.t. $r^m=r^n$

    $\implies r^m-r^n=0 \implies r^n(r^{m-n}-1)=0$

    Since $r^n \neq 0 \implies r^{m-n-1}r=1 \implies r^{-1}=r^{m-n-1} \implies R$ is a field
\end{proof}

\begin{example}
    For every prime $p$, $\Z_p$ is an integral domain, and thus a field
\end{example}
\begin{proof}
    Commutative \checkmark

    Also multiply by anything between $1$ and $p$ will not give you a multiple of $p$
    since $p$ is prime thus there are no zero divisors thus it is an integral domain
\end{proof}

\begin{definition}
    We say the characteristic of a ring $R$ is $n \in \N$ if $n$ is 
    the smallest such that for all $r\in R$ \[nr=r+r+\dots+r=0\]

    We write char$(R)=n$

    If no such $n\in\N$ exists we say char$(R)=0$
\end{definition}

\begin{proposition}
    If $R$ is a ring with 1  and char$(R)\neq 0$ then char$(R)=$ order of 1 in the groups
    $(R,+)$
\end{proposition}

\begin{proof}
    Set char$(R)=n$ and ord$(1)=n$

    $0=1+\dots+1 \implies m \mid n \implies m \leq n$

    Also $r+\dots+r\t{ (m times)}=r(1+\dots+1)=r0=0 \implies m \geq n$

    $\implies m=n$
\end{proof}

\begin{proposition}
    If $R$ is an integral domain then char$(R)=p$ or char$(R)=0$, $p$-prime
\end{proposition}

\begin{proof}
    If char$(R)=0$ we are done.

    Assume char$(R)=n \in \N$

    Assume $n=ab$ with $1\leq a,b\leq n$

    $0=1+\dots+1 \t{ (n times)}=(1+\dots+1) \t{ (a times) }(1+\dots+1)\t{ (b times)}$

    Since $R$ is an integral domain WLOG $a1=0$ 

    $\implies a \geq n \implies a=n \implies b=1$

    $\implies n$ is prime
\end{proof}

\begin{theorem}
    \label[theorem]{thm:polynomialintegral}
    $R$ is an integral domain $\iff R\left[ x \right]$ is an integral domain
\end{theorem}

\begin{proof}
    $(\implies)$
    Assume $R$ is a commutative ring with 1

    Let $f,g \in R\left[ x \right]$ s.t. 
    $f=\sum\limits^{n}_{i=0}a_ix^i$ s.t. $a_i \in R$ and 
    $g=\sum\limits^{k}_{i=0}b_ix^i$ s.t. $b_i \in R$

    WLOG let $n \geq m$

    Closure + Commutative:

    $f+g=(a_0+b_0)+(a_1+b_1)x+\dots+(a_n+b_n)x^n \in R\left[ x \right]$
    $=(b_0+a_0)+\dots+(b_n+a_n)x^n=g+f \implies (R\left[ x \right],+)$ is commutative

    \[f\cdot g= \left(\sum\limits^{n}_{i=0}a_ix^i\right)\left( \sum\limits^{k}_{i=0}b_ix^i \right)=
    \sum\limits^{m+n}_{k=0}\left( \sum\limits_{i+j=k}a_ib_j \right)x^k\]

    This is in $R\left[ x \right]$ and also since $R$ is a commutative ring 
    $\implies \sum\limits^{m+n}_{k=0}\left( \sum\limits_{i+j=k}a_ib_j \right)x^k
    =\sum\limits^{m+n}_{k=0}\left( \sum\limits_{i+j=k}b_ja_i \right)x^k = g \cdot f$

    $\implies$ multiplication is commutative in $R\left[ x \right]$

    Identity: $f+0=f=0+f$ and $f\cdot 1=f=1\cdot f$

    Inverses: Consider $-f=\sum\limits^{n}_{i=0}-a_ix^i$ $f+(-f)=0=(-f)+f$

    Associativity: $f+(g+h)=(f+g)+h$ and $f(gh)=(fg)h$ since this works for the coefficients

    Thus, $(R\left[ x \right],+)$ is an abelian group

    Distributivity: Since multiplication is commutative we only have to check one side 

    $f(g+h)=fg+fh$ since it works with coefficients

    Thus, $R\left[ x \right]$ is a commutative ring with 1

    $(\impliedby)$ Assume $R\left[ x \right]$ is a commutative ring with 1

    Let $f,g \in R \implies f,g \in R[x]$

    So all rules follow $\implies R$ is a commutative ring with 1


    $(\implies)$ Assume $R$ has zero divisors

    Let $a,b \in R$ s.t. $a,b \neq 0$ with $ab=0$

    View $a,b \in R[x] \implies ab=0 \in R[x] \implies R[x]$ has zero divisors
    
    $(\impliedby)$ Assume $R[x]$ has zero divisors

    Let $f,g \in R[x]$ s.t. $f,g \neq 0$ with $fg=0$

    Let $f=a_0+\dots+a_mx^m$ s.t. $a_m \neq 0$

    $g=b_0+\dots+b_nx^n$

    $fg=a_mb_nx^{n+m}+\dots=0 \implies a_mb_n=0 \implies R$ has zero divisors
\end{proof}